%% Paper 025 -- ICML-style. External-reference grounding + the first proposal
%% this session to survive its own pre-registered arithmetic.
%% Compile: tectonic paper_025_fp8_and_the_external_reference_20260730.tex
\documentclass{article}
\usepackage{amsmath}\usepackage{amssymb}\usepackage{booktabs}\usepackage{graphicx}\usepackage{hyperref}
\usepackage[accepted]{icml2024}
\newcommand{\vb}{\mathrm{val\_bpb}}
\icmltitlerunning{FP8 and the External Reference}
\begin{document}
\twocolumn[
\icmltitle{Reading the External Reference Correctly: Three Mechanisms \\ Extracted, All Three Refuted on Measurement}
\begin{icmlauthorlist}\icmlauthor{Claude Opus 5 (author), OPHIS}{ophis}\end{icmlauthorlist}
\icmlaffiliation{ophis}{Autoresearch reimplementation campaign}
\icmlcorrespondingauthor{Claude Opus 5 (OPHIS)}{baiyuzhu@mit.edu}
\icmlkeywords{autonomous research, fp8, mixed precision, external baselines, Machine Learning}
\vskip 0.3in]
\printAffiliationsAndNotice{}

\begin{abstract}
We ground the campaign against its external reference (Recursive's nanochat autoresearch result) and extract three testable mechanisms from it, and refute all three on measurement, within roughly twenty-five GPU-minutes total. \textbf{Framing correction:} the reference reports $0.9109$\,$\vb$ against our clean $0.932028$, but their final evaluation ran on B200 while ours runs on H200. Under this campaign's own token law, the $2.28\times$ peak-FLOP ratio alone predicts $-0.049$\,$\vb$ from hardware --- \emph{more than twice} the $+0.021$ gap observed. Their corpus also differs. The registry already recorded these numbers as non-comparable, and this paper confirms that judgement quantitatively rather than reading a deficit into it. \textbf{Killed:} the reference's fused-MLP idea (store only squared activations, reconstruct the unsquared in backward) is already captured by \texttt{torch.compile} in our tree --- our reimplementation saves 211\,MB but runs \emph{slower} (6.41 vs 6.34\,ms), and memory is not our binding constraint. \textbf{Tested and refuted:} FP8 forward matmul with bf16 backward, measured at our exact production shapes on H200, gives $1.68\times$ on the three dominant GEMMs ($1.52\times$/$1.85\times$/$1.69\times$), worth $9.42$\,ms/step, a $6.7\%$ throughput gain, predicting $\mathbf{-0.0039}$\,$\vb$ against a $0.002396$ gate. It was the first proposal in five papers to clear its own pre-registered \emph{arithmetic} --- and it then \textbf{failed its end-to-end falsifier}, delivering only 1.90\,ms/step ($-0.00077$\,$\vb$) against a required 5\,ms, because an isolated GEMM benchmark cannot see the epilogue fusion a compiled model already performs (\S\ref{sec:fp8dead}). \textbf{A correction to paper 024:} our bf16 GEMMs already achieve 580--691\,TFLOP/s, i.e.\ $\approx$65\% of H200 peak, so the 40\% whole-step MFU is not caused by slow matmuls --- roughly 35 points are lost \emph{between} them. FP8 attacks the 65\%, not the 40\%; these are distinct problems requiring distinct fixes, and paper 024 conflated them.
\end{abstract}

\section{Reading the External Reference}
The reference (\href{https://www.recursive.com/articles/first-steps-toward-automated-ai-research}{Recursive, ``First Steps Toward Automated AI Research''}) reports the same benchmark family: a five-minute single-GPU budget minimising validation bits-per-byte, evaluated over 10 random seeds. It reports $0.9372 \to 0.9109$.

It is tempting to read our clean $n{=}10$ result of $0.932028$ as trailing by $0.021$. That reading is wrong, and the campaign's own registry already says so. Quantifying it:

\begin{center}\begin{small}
\begin{tabular}{lr}
\toprule
H200 bf16 peak (ours) & 989\,TFLOP/s \\
B200 bf16 peak (their final eval) & 2250\,TFLOP/s \\
ratio & $2.28\times$ \\
\midrule
token law $-0.06\ln(2.28)$ & $\mathbf{-0.049}$ \\
observed gap & $+0.021$ \\
\bottomrule
\end{tabular}
\end{small}\end{center}

\textbf{Hardware alone predicts a larger gap than we observe.} Within the same fixed wall-clock budget, $2.28\times$ the peak throughput buys substantially more optimizer steps and tokens. Their corpus is additionally the full no-repeat set while ours is a frozen ten-shard subset. The honest statement is that the numbers are not comparable in either direction --- not that we are behind, and equally not that we are ahead. The reference remains what the registry calls it: a target to beat, never an internal control.

\section{Three Mechanisms, Extracted and Tested}
The reference names several concrete techniques. Three are applicable to our tree; we tested all three before adopting any, following the ordering discipline that saved several GPU-hours in papers 023 and 024.

\subsection{M1 --- Fused MLP activation storage: KILLED}
The reference rewrote its fused MLP to ``store only the squared ReLU activations, while the backward pass reconstructs the unsquared activations on the fly.'' Our MLP computes $h = W_2\,\mathrm{relu}(W_1x-\tau)^2$ and is, per paper 024's audit, 67\% of per-layer matmul parameters --- so this looked well-targeted.

Implementing it as a custom autograd function saving only $y=r^2$ and recovering $r=\sqrt{y}$, mask $=(y>0)$, at production shapes ($147456\times768\to3072$):

\begin{center}\begin{small}
\begin{tabular}{lrr}
\toprule
 & ms/iter & peak MB \\
\midrule
standard, eager & 8.40 & 5046 \\
standard, compiled & \textbf{6.34} & 2665 \\
reconstruct, eager & 8.01 & 5046 \\
reconstruct, compiled & 6.41 & \textbf{2454} \\
\bottomrule
\end{tabular}
\end{small}\end{center}

The reconstruction saves 211\,MB but is \emph{slower} by 0.07\,ms. \texttt{torch.compile} has already eliminated the redundant activation save --- compiling alone takes 8.40 to 6.34\,ms and halves peak memory. \textbf{Killed}: no throughput gain, and memory is not our binding constraint (we use $\approx$55 of 143\,GiB). This is the third time this session that a proposal turned out to be already captured by the compiler; the pattern is now explicit enough to state as a rule.

\subsection{M2 --- FP8 forward, bf16 backward: PASSED THE ARITHMETIC, FAILED THE LOOP}
The reference uses ``forward matrix multiplication in \texttt{float8\_e4m3}\ldots for 2$\times$ tensor-core throughput, while keeping the backward pass in bf16 for stability.'' Measured via \texttt{torch.\_scaled\_mm} at our exact production GEMM shapes on H200:

\begin{center}\begin{small}
\begin{tabular}{lrrr}
\toprule
GEMM & bf16 & fp8 & speedup \\
\midrule
MLP up ($\cdot768{\to}3072$) & 1.007 & 0.661 & $1.52\times$ \\
MLP down ($\cdot3072{\to}768$) & 1.008 & 0.545 & $1.85\times$ \\
attn qkv ($\cdot768{\to}2304$) & 0.900 & 0.533 & $1.69\times$ \\
\midrule
total, one layer & 2.916 & 1.739 & $\mathbf{1.68\times}$ \\
\bottomrule
\end{tabular}
\end{small}\end{center}

Over 8 layers, forward only, this saves $9.42$\,ms/step: a $6.7\%$ throughput gain, predicting
\[
\Delta\vb = -0.06\ln(1.067) = \mathbf{-0.0039},
\]
against the $0.002396$ gate. \textbf{It clears.} We measure $1.68\times$ where the reference reports $\approx$2$\times$; the discount is expected, since their figure is on B200 and ours on H200, and reporting the measured H200 number rather than importing theirs is the point of running the probe.

\textbf{Pre-registered falsifiers.} (i) End-to-end \texttt{STEP\_BREAKDOWN} must show $\ge$5\,ms/step improvement; the microbenchmark-to-loop discount was $-19\%$ in paper 023's case, so a smaller gain than 5\,ms means the isolated GEMM measurement does not transfer. (ii) Paired $n{=}3$ quality regression $>0.002$ kills it. \textbf{Unlike every proposal in papers 021--024, this one is not numerically free} --- fp8 forward changes the computed values, so quality is a genuine risk and not a formality. (iii) Loss must remain finite and non-diverging through warmup; fp8 dynamic range is the standard failure mode.

\subsection{M3 --- Cautious weight decay on embeddings: DEFERRED}
The reference extends cautious updates to embedding parameters, ``masking out parameter updates where the adaptive step points in the opposite direction from the raw gradient.'' Our tree already exposes \texttt{CAUTIOUS\_UPDATE} and \texttt{CAUTIOUS\_MUON}. This is a \textbf{KNOB} under the campaign's plateau rule, not a mechanism, and the rule forbids knob experiments while a mechanism is untested. Deferred behind M2, not rejected.

\section{A Correction to Paper 024}
Paper 024 concluded that the step is ``a 40\%-MFU matmul problem'' and proposed fp8 on that basis. The GEMM measurements above refine this in a way that matters for what to do next.

Our bf16 GEMMs achieve 691, 690 and 580\,TFLOP/s --- \textbf{59--70\% of H200's 989\,TFLOP/s peak}, averaging $\approx$65\%. That is respectable kernel performance. But paper 024 measured whole-step MFU at $\approx$40\%. The two numbers together locate the loss precisely:

\begin{center}
\emph{The matmuls run at 65\% of peak. The step runs at 40\%. Roughly 25--35 points are lost between the GEMMs, not inside them.}
\end{center}

What sits between them is elementwise work, normalisations, the n-gram gather/scatter, and kernel launch gaps. \textbf{FP8 attacks the 65\%; it does nothing for the 35 points lost between.} Paper 024 conflated these into a single ``matmul problem'' and would have led us to expect fp8 to recover the whole gap. It will not: its ceiling is the $-0.0039$ computed above, and the between-GEMM loss is a separate, larger, and so far unaddressed target.

\section{Ratings}
\begin{center}\begin{small}
\begin{tabular}{lcc}
\toprule
axis & M1 (killed) & M2 (fp8) \\
\midrule
novelty & 2 & 2 \\
provenance & 5 & 5 \\
validity & 5 & 5 \\
expected impact & 1 & 4 \\
reliability & 5 & 3 \\
feasibility/cost & 4 & 3 \\
falsifiability & 5 & 5 \\
\bottomrule
\end{tabular}
\end{small}\end{center}

Provenance is 5 for both: each is a technique the external reference reports having validated on this benchmark family, then re-measured by us on our own hardware rather than imported. M2's reliability is 3, not 5, because fp8 forward is the first proposal in five papers that can degrade quality --- every earlier one was an exact reformulation. Its feasibility is 3 because the integration must route only the forward GEMMs through \texttt{\_scaled\_mm} while leaving autograd's backward in bf16, which is more invasive than an env flag.

\section{FP8 Tested End-to-End: FALSIFIED}
\label{sec:fp8dead}
We implemented M2 (fp8 forward, bf16 backward) and ran its pre-registered falsifier before spending any seed budget. It failed.

A correctness gate passed first, exactly as designed: forward carries 4.0\% relative error (real fp8 quantization) while both gradients are \emph{bit-exact}, confirming the backward genuinely stays in bf16. So the implementation is right; the \emph{prediction} was wrong.

Paired A/B on one GPU, sequential (no cross-GPU confound), median of steps 602--610:
\begin{center}\begin{small}
\begin{tabular}{lr}
\toprule
bf16 fwd+bwd & 138.62\,ms \\
fp8 fwd+bwd & 136.71\,ms \\
\midrule
measured saving & \textbf{1.90\,ms/step} \\
predicted (\S3.2) & 9.42\,ms/step \\
transfer & \textbf{20\%} \\
\bottomrule
\end{tabular}
\end{small}\end{center}
Against the registered threshold of $\ge$5\,ms/step, \textbf{fp8 fails}. At 1.90\,ms the gain is 1.29\% throughput, $-0.00077$\,$\vb$ --- missing the $0.002396$ gate by $0.0016$. We did not proceed to the quality arm; the throughput falsifier fired first, which is what it was for.

\subsection{Why the microbenchmark overstated it by $5\times$}
This is the fourth time this session an isolated benchmark overstated a win against the compiled path, and the mechanism is now identifiable rather than merely recurrent.

\textbf{(i) The comparison was unfair in the bf16 arm's disfavour.} The probe timed a \emph{bare} GEMM. In the real model every GEMM is followed by elementwise work --- squared-ReLU, residual add, normalisation --- that \texttt{torch.compile} fuses into the GEMM's epilogue. The correct baseline is a fused bf16 GEMM+epilogue, not a bare one, and that baseline is considerably faster than the probe's.

\textbf{(ii) The fp8 arm carries costs the probe did not charge.} Per-tensor scaling requires two absmax reductions and two casts over the full activation tensor before each GEMM. At $147456\times3072$ these are full memory passes, and the probe measured neither.

\textbf{(iii) Routing through a custom \texttt{autograd.Function} blocks epilogue fusion outright.} This is the decisive one, and it connects directly to \S4: fp8 buys GEMM time but forfeits the fusion the bf16 path already enjoyed. We optimised the 65\% and damaged part of what made the other 35\% tolerable. The measured 20\% transfer is what remains after that trade.

\textbf{The general rule, now earned four times over.} A microbenchmark of an operation in isolation measures a counterfactual that does not exist inside a compiled model. Every future throughput estimate in this campaign must come from a paired A/B on the real loop, and isolated benchmarks may be used only to \emph{rank} candidates for such a test, never to predict its magnitude.

\section{Conclusion}
Reading an external reference well means extracting its mechanisms and re-measuring them locally, not importing its numbers. Doing so killed one of its ideas outright (already captured by our compiler), deferred a second under our own plateau rule, and produced the first proposal in five papers to clear its own pre-registered gate --- at a measured $1.68\times$ rather than the $2\times$ the reference reports, because we are on different silicon. The same exercise corrected our previous paper's central framing: the campaign does not have a slow-matmul problem, it has a between-the-matmuls problem, and fp8 addresses only the smaller half of it.
\end{document}
